\section{Trabalhos Relacionados}
\label{sec:related}

\subsection{\textit{Datasets} de toxicidade em texto}

A produção de \textit{datasets} rotulados para detecção de toxicidade tem
sido dominada por iniciativas em língua inglesa.

\paragraph{Wikipedia Talk Corpus / Jigsaw.}
Wulczyn~et~al.~\cite{wulczyn2017} apresentaram o \textit{Wikipedia Talk
Corpus}, base posteriormente expandida para o
\textit{Jigsaw}~\cite{borkan2019nuanced} com 160 mil comentários e seis
categorias binárias.
\textbf{Pontos fortes:} grande escala, anotação multi-rótulo, ampla adoção
(\textit{benchmark} \textit{de facto} para classificação de toxicidade em
inglês). \textbf{Pontos fracos:} (i)~rótulos binários descartam a
incerteza inter-anotadores; (ii)~não há justificativa amostral dos
\textit{splits} treino/teste; (iii)~viés de auditoria por
\textit{Wikipedians} não representa uso público de mídias sociais.
\textbf{Razão de não-adoção:} idioma inglês incompatível com o objetivo
deste trabalho.

\paragraph{Civil Comments com atributos de identidade.}
Borkan~et~al.~\cite{borkan2019nuanced} estenderam o Jigsaw com 1{,}8
milhões de comentários e atributos de identidade (gênero, raça, religião,
orientação sexual). \textbf{Pontos fortes:} cobertura demográfica fina e
métricas \textit{subgroup-aware}. \textbf{Pontos fracos:} sub-rotulagem
extrema dos atributos de identidade --- a maior parte dos comentários não
recebe nenhum atributo, comprometendo a estimação de viés por
\textit{subgroup}. \textbf{Razão de não-adoção:} idioma e ausência de cauda
estratificada explicitamente comparável.

\paragraph{Datasets de \textit{tweets}.}
Davidson~et~al.~\cite{davidson2017} e Founta~et~al.~\cite{founta2018}
construíram \textit{datasets} sobre \textit{Twitter} com foco em discurso de
ódio explícito e abusivo. \textbf{Pontos fortes:} taxonomias detalhadas e
volume relevante. \textbf{Pontos fracos:} (i)~classes definidas operacionalmente
(\textit{hate} vs \textit{offensive} vs \textit{neither}) com baixa
confiabilidade inter-anotadores; (ii)~discutida fragilidade contra
\textit{tweets} \textit{adversariais}~\cite{davidson2017}.
\textbf{Razão de não-adoção:} idioma inglês.

\paragraph{HateXplain.}
Mathew~et~al.~\cite{mathew2021} introduziram um \textit{dataset} de 20 mil
amostras com \textit{rationales} explicáveis para o rótulo. \textbf{Pontos
fortes:} balanceamento artificial de classes (cada classe $\approx 33\%$),
o que viabiliza alocação proporcional sem caudas raras. \textbf{Pontos fracos:}
o balanceamento construído \textit{a priori} elimina o fenômeno de
desbalanceamento extremo --- precisamente o que motiva a alocação ótima
em corpora reais. \textbf{Razão de não-adoção:} sua estrutura balanceada o
torna inadequado como caso de teste para Neyman; resolveria nosso problema
por construção.

\paragraph{Contexto conversacional.}
Pavlopoulos~et~al.~\cite{pavlopoulos2020} mostraram que adicionar contexto
conversacional altera os rótulos em uma fração não desprezível dos casos.
\textbf{Implicação:} qualquer auditoria amostral precisa fixar o protocolo
de rotulagem (com ou sem contexto), o que reforça a importância de
documentar o esquema do \textit{corpus}-fonte.

\paragraph{Datasets em português brasileiro.}
Leite~et~al.~\cite{leite2020,leite2020_thesis} publicaram o \textbf{ToLD-Br}
(seis categorias, votos $0$--$3$). \textbf{Pontos fortes:} esquema ordinal de
votos preserva incerteza; volume relativamente grande ($N=21$ mil) para um
\textit{corpus} em português; cobre o espectro \textit{LGBTfobia–racismo–
misoginia–xenofobia}. \textbf{Pontos fracos:} sem justificativa amostral no
artigo original; estratificação multi-rótulo não-trivial.
\textbf{Razão de adoção:} fit ideal para o presente estudo (idioma, escala,
multi-rótulo, caudas raras explícitas).

Vargas~et~al.~\cite{vargas2022hatebr,vargas2022_thesis} apresentaram o
\textbf{HateBR}, com 7 mil comentários do \textit{Instagram} e rotulagem
hierárquica binária (ofensivo vs não, com sub-categorias).
\textbf{Pontos fortes:} qualidade de anotação por especialistas;
contextualização cultural brasileira distinta. \textbf{Pontos fracos:}
$N$ pequeno limita estratos raros (xenofobia, capacitismo); rotulagem
hierárquica não permite comparação direta com ToLD-Br.
\textbf{Razão de não-adoção como fonte primária:} escala insuficiente para
ilustrar Neyman vs proporcional --- avaliado como \textit{corpus} de
extensão futura.

Trajano~et~al.~\cite{trajano2024olid} disponibilizaram o
\textbf{OLID-BR} ($\approx 7$ mil exemplos) com taxonomia OLID adaptada.
\textbf{Pontos fortes:} alinhamento com a tradição internacional OLID,
facilitando comparação \textit{cross-lingual}. \textbf{Pontos fracos:}
escala similar à do HateBR; sem esquema ordinal de votos.
\textbf{Razão de não-adoção:} mesmas considerações do HateBR.

\subsection{Auditoria de modelos de linguagem e viés}

\paragraph{RealToxicityPrompts.}
Gehman~et~al.~\cite{gehman2020} popularizaram o uso de \textit{prompts}
extraídos de \textit{datasets} de toxicidade como instrumento de auditoria
de modelos generativos. \textbf{Pontos fortes:} protocolo reproduzível e
amplamente adotado pela comunidade. \textbf{Pontos fracos:} a amostragem dos
\textit{prompts} é por escore (\textit{toxic} vs \textit{non-toxic}) sem
justificativa amostral formal --- a quantidade de \textit{prompts} é
arbitrária. \textbf{Posicionamento:} este artigo cobre exatamente esta
lacuna inferencial, embora aplicada a auditoria de \textit{datasets} e não
ao protocolo de \textit{prompting}.

\paragraph{Métricas \textit{subgroup-aware}.}
Borkan~et~al.~\cite{borkan2019nuanced} formalizaram métricas como
\textit{Equality of Opportunity} e \textit{Demographic Parity} para medir
viés não-intencional. \textbf{Ponto crítico:} todas dependem de tamanhos
amostrais por subgrupo, mas raramente acompanhadas de derivação amostral
formal --- de novo, lacuna que este trabalho preenche.

\paragraph{Crítica metodológica em escala.}
Bender~et~al.~\cite{bender2021} sustentam que a auditoria de modelos de
linguagem em larga escala precisa de fundamentos estatísticos explícitos,
para além de \textit{benchmarks} pontuais. \textbf{Posicionamento:} este
artigo é uma resposta concreta a esse chamado.

\subsection{Aplicações de teoria amostral em PNL}

A teoria amostral
clássica~\cite{cochran1977,bolfarine2005,lohr1999,thompson2002,cordeiro2010}
é prática consolidada em ciências sociais, epidemiologia e pesquisa de
mercado, mas é pouco utilizada em PNL aplicada. Os \textit{datasets} citados
acima geralmente apresentam apenas a divisão treino/validação/teste sem
justificar estatisticamente a representatividade dos
\textit{splits}~\cite{leite2020,vargas2022hatebr,davidson2017}.
Quando a amostragem é discutida, o foco é geralmente \textit{class balancing}
para treino preditivo --- e não estimação inferencial de parâmetros
populacionais. \textbf{Lacuna identificada:} ausência de aplicação explícita
de alocação ótima e estimadores tipo razão e regressão a \textit{corpora} de
auditoria de toxicidade.

\subsection{Posicionamento deste trabalho}

A Tabela~\ref{tab:related-comparison} sumariza a comparação. Os trabalhos
discutidos cobrem (a) a construção de \textit{datasets} de toxicidade em
diferentes línguas e (b) métricas de auditoria sobre modelos. Persiste,
entretanto, a lacuna metodológica de \textit{aplicar amostragem
probabilística estratificada} sobre tais \textit{corpora} para estimar
parâmetros populacionais com erro-padrão controlado. O presente artigo se
posiciona neste espaço, aplicando alocação de Neyman e estimadores tipo
razão e regressão de forma reproduzível sobre o ToLD-Br.

\begin{table}[!t]
\centering
\small
\caption{Comparação resumida das abordagens dos trabalhos relacionados.}
\label{tab:related-comparison}
\begin{tabular}{lcccc}
\toprule
Trabalho & Idioma & Justifica $n$? & Trata caudas raras? & Estima $\bar Y$? \\
\midrule
\cite{wulczyn2017}        & EN & não & não & não \\
\cite{borkan2019nuanced}  & EN & não & sim (identidade) & não \\
\cite{leite2020}          & PT & não & não & não \\
\cite{vargas2022hatebr}   & PT & não & sim (hierárquico) & não \\
\textbf{Este trabalho}    & \textbf{PT} & \textbf{Cochran} & \textbf{Neyman+censo} & \textbf{razão+regressão} \\
\bottomrule
\end{tabular}
\end{table}
